\documentclass[11pt]{article}
\usepackage[utf8]{inputenc}
\usepackage[T1]{fontenc}
\usepackage{lmodern}
\usepackage[margin=1in]{geometry}
\usepackage{amsmath,amssymb}
\usepackage{textcomp}
\usepackage[dvipsnames]{xcolor}
\usepackage{enumitem}
\usepackage{parskip}
\usepackage{titlesec}
\usepackage[round]{natbib}
\usepackage[colorlinks=true,linkcolor=blue,citecolor=blue,urlcolor=blue]{hyperref}

\titleformat{\section}{\large\bfseries}{\thesection}{1em}{}
\titleformat{\subsection}{\normalsize\bfseries}{\thesubsection}{1em}{}
\titleformat{\subsubsection}{\normalsize\bfseries\itshape}{}{0em}{}

\title{\textbf{Certified Sound Alone, Silently Unsound Together:\\ Measuring Emergent Shortcuts under Boolean Verifier Composition}}
\author{}
\date{}

\newcommand{\ph}{\textit{[placeholder --- fill in after experiments are run]}}

\begin{document}
\maketitle
\vspace{-2em}

\begin{quote}
\small\textit{Assembled from the six confirmed stage outputs of an \texttt{auto-research-discussion} run (seed topic B1.4), post research-taste audit (\textcolor{ForestGreen}{\textbf{PASS}}). Status: \textbf{specification / paper draft --- no experiments have been run.} All numbers below are pre-registered design targets, not results.}
\end{quote}

\begin{abstract}
Verification pipelines for LLM agents and code generation increasingly gate acceptance on stacks of individually-hardened checks --- a test suite plus a style linter, an outcome verifier plus a flawed-reasoning detector --- under the implicit assumption that combining sound checks can only tighten the gate. This survives an obvious counterargument only by accident: the Boolean composition of two acceptors sound on a fixed task in the abstract is trivially sound, but no deployed verifier is certified that way --- each is sound only relative to the facet it checks, its own task variant, and a finite search budget, and it is this relativized notion that composition puts at risk. We construct matched task variants whose single-verifier soundness is exhaustively, non-vacuously certified before composition, then search under identical adversarial budgets for a candidate a composed AND/OR gate accepts while the ground truth rejects. A mandatory re-attack of every ``emergent'' shortcut's single-gate origin at ten times its certification budget rules out unequal search effort, and a structural taxonomy regressed on a coverage-complementarity score separates genuine composition effects from verifiers already blind to the same region. We report whether, and under what conditions, composing two certified-sound verifiers manufactures a shortcut neither admits alone --- anchoring the headline claims not in the constructed pilot, whose existence result is conditional on a pre-registered candidate-space extension rule and serves chiefly to validate the instrument, but in the measured shortcut rate on realistic benchmark stacks, where no constructed extension exists, and in whether a complementarity signal computable before any attack predicts where that rate concentrates.
\end{abstract}

\section{Introduction}

Verification is rarely a single check anymore. As large language models are put in charge of generating code, executing agentic plans, and producing content that must clear a bar before it is accepted, rewarded, or deployed, the checks guarding that bar have started to stack. Reinforcement-learning pipelines built on verifiable rewards decompose a single reward signal into a chain of content and style checks; outcome-only verifiers get a second, process-level check bolted on to catch correct-answer-wrong-reasoning cases; execution-based test suites for software-engineering agents are paired with execution-free reward models to cover what tests miss; safety guardrails aggregate a whole panel of per-category classifiers into one verdict. Every one of these designs shares an unstated premise: that adding a second, independently defensible check can only make the gate stricter, never looser. It is a natural premise, and it is exactly the kind of premise that is worth checking before it is relied on any further.

The trouble is that the obvious argument for this premise is both true and irrelevant. Extensionally, on one fixed task, the conjunction of two sound acceptors is sound --- it accepts a subset of what either accepts alone --- and their disjunction is sound in the complementary sense of never being stricter than either. That fact is not in dispute. What is in dispute is whether it describes the situation any real pipeline is actually in. No deployed verifier is ever certified sound in that abstract, task-global sense; it is certified sound only relative to the specific facet it checks (a unit test checks input-output behavior, not resource usage; a style linter checks formatting, not correctness), its own task variant, and a finite budget of adversarial or exhaustive search. Boolean composition changes the distribution of candidates a solver will actually produce --- an AND-gate's second constraint can push solutions into a region the first check's certificate was never tested against; an OR-gate's weaker branch can admit exactly what the stronger branch was built to reject. The premise that composition only tightens the gate is true for acceptors nobody actually builds, and unexamined for the ones everybody does.

The nearest existing work approaches this gap from adjacent but distinct angles, and none of it closes it directly. Systems that already compose checks validate the composition only by its downstream effect, never by examining the combined gate's own soundness: a reward-chain framework that converts one verifiable answer into a hierarchy of content and style sub-rewards \citep{rewardchain2601} motivates its design by the hacking it avoids, and a flawed-positive-aware policy optimizer adds a process-level detector to an outcome verifier \citep{fapo2510} and reports the resulting policy's end-task gains --- in neither case is the AND-composed gate itself ever attacked to see whether it, as a whole, still holds; notably, the reward-chain framework's own composition is a soft weighted average of sub-rewards rather than a hard Boolean gate, so even its downstream-only validation does not directly speak to Boolean AND/OR composition. A guardrail built on reasoning over per-category safety scores \citep{r2guard2407} documents, in its own motivating analysis, that naively aggregated classifiers already under-estimate risk when unsafe intent is spread across several weak signals --- real, cited evidence that composition can leak --- but the proposed fix targets homogeneous categories inside one guardrail, not the heterogeneous pairs (a test suite and a judge, an execution check and a reward model) this paper is concerned with. A formal account of reward hacking for a single proxy \citep{correlatedproxies2403} formalizes when a proxy is valid relative to a reference policy, a single-proxy theory that never engages the discrete accept/reject algebra of a hard gate or says anything about how such validity composes across two checks. And on the other side, work that hardens or maps a verifier's failure surface --- adversarially strengthening a test suite against the current solver \citep{utrl2508}, or mining a judge's own blind spots from its misjudgments \citep{biasscope2602} --- does so one verifier at a time, with no account of what happens when two independently mapped failure surfaces are combined.

Turning the gap into something measurable surfaces four distinct obstacles. First, the phenomenon has no content unless it is made non-vacuous: any claimed ``emergent'' shortcut has to be built against task variants where each verifier's soundness is genuinely, checkably certified beforehand, not casually assumed, since the abstract composition argument is trivially true and would otherwise swallow any counterexample as an artifact of an unstated assumption. Second, one hand-built counterexample proves only that the phenomenon is possible, not how often it recurs --- turning existence into a rate requires a separately designed, budget-matched sampling protocol across a realistic grid of tasks and verifier pairings, since the space of pairings and shortcut strategies grows too combinatorially to explore by constructing more examples by hand. Third, because certification outside a fully enumerated pilot is itself only an absence-of-evidence claim bounded by a search budget, any discovered composed-gate shortcut has to be checked against the possibility that the attacker simply spent more effective effort against the composed gate than the single gates were certified with --- a budget-symmetry confound distinct from any question about what the verifiers actually cover. Fourth, and separately from budget symmetry, a shortcut that both verifiers would have missed independently even without composition is a different finding from one that composition itself brought into being, and no single statistic distinguishes the two without first characterizing each verifier's failure surface on its own.

Our approach turns each obstacle into a specific design decision. A bounded, exhaustively enumerable candidate space lets per-verifier soundness be certified as a machine-checked fact rather than an assumption, closing off the vacuity trap before any composed gate is built. An identical, budget-matched adversarial search is run against every single and composed gate across a realistic task and verifier-pairing grid, converting a hand-built existence result into an estimable rate rather than a curated anecdote. Any shortcut provisionally labeled ``emergent'' triggers a mandatory re-attack of its single-gate origin at ten times the original certification budget, so a finding survives only if it cannot be explained away as an under-searched single verifier. And every surviving shortcut is classified by an ablation triple against each single gate and the composed gate, sorted into a small structural taxonomy, and related to a measurable coverage-complementarity score computed directly from the two verifiers' blind spots --- turning ``these two checks happened to be complementary'' from an after-the-fact story into a quantity that can be checked in advance of running any attack at all.

We evaluate this design on a hand-constructed pilot suite where per-verifier soundness is exhaustively decidable, then scale to a realistic grid spanning general-purpose code-synthesis benchmarks and benchmarks with natively multi-constraint tasks, comparing naive Boolean composition against its single-verifier components, an adversarially hardened single verifier, and a structured (non-naive) aggregation scheme.

In summary, we contribute:
\begin{itemize}[leftmargin=1.4em]
\item A non-vacuous, machine-checkable construction of matched task variants that makes facet-relative soundness a provable property rather than an assumption, closing off the trivial-soundness objection before any composed-gate claim is made.
\item A budget-stamped certification and mandatory re-attack protocol that lets a composed-gate shortcut be attributed to composition itself, ruling out unequal search effort as an alternative explanation.
\item A matched-budget adversarial search design that converts a constructed existence result into an estimable prevalence rate across a realistic benchmark and verifier-pairing grid, and directly compares hardening one verifier against adding a second.
\item A structural taxonomy and a coverage-complementarity signal, computable before any adversarial search is run, that separates genuine composition effects from verifiers that were already blind to the same region and predicts where composition-induced gaps concentrate.
\end{itemize}

\section{Problem Definition}

\noindent\textbf{Formal setup.} Consider a task $T = (\text{spec}, r, D)$ with reference solution $r$ and differential-test corpus $D$, $|D| \geq 10^4$, and a checkable ground truth $GT(c) = 1$ iff candidate $c$ agrees with $r$ on every input in $D$. A verifier $V$ maps a candidate to $\{\text{accept}, \text{reject}\}$. A gate $G$ is either a single verifier or a Boolean composition, $G \in \{A, B, A \land B, A \lor B\}$, of two verifiers $A$ and $B$. Task variants are a tuple $\langle T, t_A, t_B, t_\land, t_\lor \rangle$ sharing spec and full-task ground truth, but each variant $v$ carries its \emph{own} pre-registered certification space $S_v \subseteq S_\circ$: the composed variants' space $S_\circ$ extends each single variant's space $S_A$ / $S_B$ by exactly the constructs the added gate makes reachable, per an extension rule fixed before any certification runs. This is what makes ``the candidate distribution its own gate induces'' operational: each verifier is certified over the space its own gate's variant induces, never over the composed space --- a degenerate configuration where $S_A = S_B = S_\circ$ and both verifiers are certified against one shared space is provably incapable of admitting an emergent shortcut (Section~\ref{sec:baselines-placebo} below) and is reserved as the null-by-construction placebo rather than treated as a general-case setup. A \emph{shortcut} against a gate $G$ on a task variant $t$ is a candidate $c$ with $G(c) = \text{accept}$ and $GT(c) = 0$. Soundness is defined relative to a certification procedure $\Pi$: $\mathrm{Sound}(V, t, \Pi) = 1$ iff $\Pi$ finds no shortcut against $V$ on $t$'s own space $S_t$. $\Pi$ takes one of two forms with different epistemic status --- $\Pi = \text{exhaustive}$, a complete decision procedure over the bounded space $S_t$ (used in the pilot, where soundness claims are proofs), or $\Pi = \text{budgeted}(B_{\text{cert}})$, an adversarial search by a dedicated certification attacker $\mathcal{A}_{\text{cert}}$ --- a different model, prompt scaffold, and seed family from the comparison attacker $\mathcal{A}$ used in the experiments below, so that certification and attack draw independent randomness --- stopped after a stated budget with zero unsound acceptances (used at scale, where a soundness claim is an absence-of-evidence statement relative to $B_{\text{cert}}$, not a certificate that a larger budget could not overturn). The object of study is whether there exists a task family and a composition rule $\circ \in \{\land, \lor\}$ such that $\mathrm{Sound}(A, t_A, \Pi) = \mathrm{Sound}(B, t_B, \Pi) = 1$ over their own spaces $S_A, S_B$ while the composed variant $t_\circ$ admits a shortcut against $A \circ B$ over $S_\circ$ --- and, where this occurs, under what conditions (formalized below as a coverage-complementarity score) it recurs.

\noindent\textbf{Why this is hard.} The composition of two verifiers that are sound in the fully abstract, task-global sense is provably sound --- this is not the claim under test. What makes the question non-trivial is that $\mathrm{Sound}(V, t, \Pi)$ is relative to the facet $V$ actually checks, the task variant $t$ it was certified on, and, outside the exhaustive pilot, a finite search budget rather than a complete proof --- so a construction has to make each single-verifier certificate genuinely checkable before any composed variant is built, or an ``emergent'' finding would just be an artifact of an unstated assumption rather than a real result. Once that non-vacuous construction exists, a single instance only demonstrates possibility: estimating a rate over realistic task and verifier-pairing combinations requires a separately designed, budget-matched sampling protocol, since the combinatorics of pairings and shortcut strategies rule out brute-force enumeration beyond the pilot. Any shortcut discovered against a composed gate must then survive two distinct attribution checks before it can be called a genuine effect of composition: first, that the attacker was not simply given more effective search effort against the composed gate than the single-verifier certificates were earned with (a budget-symmetry confound), and second --- a separate concern that persists even under perfectly matched effort --- that the shortcut is not something both verifiers would have missed independently, a shared blind spot with no composition-specific mechanism behind it at all.

\medskip
\noindent\textbf{What each stage can and cannot show.} One consequence of the per-variant space construction is scoped explicitly here, because it determines the evidential role of each stage. In the pilot, the composed variant's space extends the single variants' spaces by a pre-registered rule --- so a discovered emergent shortcut there demonstrates that facet-relative certification is not closed under gate-induced growth of the reachable candidate space, a property of the certification regime exhibited by construction, not a surprising algebraic fact about Boolean operators. The pilot's job is therefore to validate the instrument end-to-end and pin the definitions down as machine-checked objects; it is not the paper's headline evidence. The load-bearing empirical claims belong to the scale-up, where candidates are free-form programs against real benchmark tasks, no constructed extension rule exists, and what is measured --- the rate at which a budget-matched adversary finds composed-gate shortcuts that survive attribution, and whether a pre-computable complementarity signal predicts where they concentrate --- is a property of verifier stacks as actually deployed.

\section{Method Overview}

At a conceptual level, the study is one measurement pipeline built around a single non-vacuity commitment: nothing is asserted about a composed gate until the pieces it is composed from have their own soundness demonstrated, not assumed. Task variants and verifier pairs enter the pipeline already carrying machine-checked single-verifier certificates; an identical adversarial search is then pointed at every single and composed configuration under matched budgets; any candidate the composed gate accepts but the ground truth rejects triggers an attribution check before it counts as anything at all; and only shortcuts that survive that check are handed to a classifier that separates genuine composition effects from verifiers that were already blind to the same region and relates the outcome to a measurable property of the verifier pair itself.

The \textbf{bounded-DSL task-pair generator} builds, for a shared task and ground truth, the single-A, single-B, AND, and OR variants over a candidate space small enough to be exhaustively enumerated in the pilot, so per-verifier soundness starts as a proof rather than an assumption. The \textbf{budget-stamped facet certifier} attaches to every single-verifier certificate the exact procedure and budget it was earned under, and mandates a $10\times$-budget re-attack of any provisionally ``emergent'' shortcut's single-gate origin before the label is finalized, closing the budget-symmetry gap. The \textbf{budget-matched adversarial attacker} applies one fixed attack configuration identically to every single and composed gate across a realistic benchmark and verifier-pairing grid, turning a hand-built existence result into an estimable rate and enabling a direct comparison between hardening one verifier and adding a second. The \textbf{composition-gap classifier} takes every shortcut that survives the re-attack check, labels it with an ablation triple against each single gate and the composed gate, sorts it into a structural taxonomy, and regresses its incidence on a coverage-complementarity score computed from the two verifiers' independently mapped blind spots.

\section{Methodology}

\subsection{Bounded-DSL Task-Pair Generator}
\textbf{Challenge.} The Boolean composition of two verifiers that are sound in the abstract, task-global sense is provably sound, so any claimed exception has no content unless each single-verifier certificate is demonstrably non-vacuous --- genuinely earned, not assumed --- before the composed variant is even constructed.

\textbf{Insight.} The vacuity objection cannot be argued away; it can only be closed constructively, by building a candidate space small enough that ``no shortcut passes this verifier alone'' is a machine-checked fact rather than a claim taken on faith.

\textbf{Method.} For a shared task and a fixed reference solution plus a differential-test corpus of at least $10^4$ inputs, the generator constructs a bounded program representation --- a fixed abstract-syntax-tree size and opcode budget --- small enough to be exhaustively enumerable in the pilot, and produces four matched variants sharing that ground truth: a single-verifier-A variant, a single-verifier-B variant, a conjunctive (AND) variant, and a disjunctive (OR) variant. Each variant carries its own pre-registered certification space rather than one space shared across all four: the composed variants' space extends each single variant's space by exactly the constructs the added gate makes reachable (e.g., a resource constraint that admits lookup-table forms absent from the single-verifier-A space), per an extension rule fixed in the generator's manifest before any certification runs. Exhaustive enumeration of each variant's own bounded space is used to certify the corresponding single-verifier variant before the composed variant, over its own (larger) space, is examined --- certifying over one shared space instead would make the composed-gate question decidable in advance by pure Boolean algebra (Problem Definition) and is used only for the degenerate placebo configuration, never as the general construction.

\subsection{Budget-Stamped Facet Certifier}
\textbf{Challenge.} Outside the exhaustively enumerable pilot, a soundness certificate is an absence-of-evidence claim relative to a search budget, not a proof --- so a composed-gate shortcut discovered at scale could be explained by unequal search effort rather than by composition, and nothing in a one-shot certificate rules that out.

\textbf{Insight.} The only way to make a budget-relative certificate mean something is to actively try to break it: re-spending disproportionate search effort specifically against the single-gate variants of any provisionally ``emergent'' shortcut checks whether it was really just an under-budgeted single-verifier miss.

\textbf{Method.} Every single-verifier certificate is stamped with the exact certification procedure and budget under which it was produced --- full enumeration in the pilot, or a stated adversarial attack budget at scale. Any shortcut that appears to pass a composed gate while both single-verifier variants remain certified sound at the matched budget is provisionally labeled emergent; before that label is finalized, the single-gate variants are re-attacked at ten times the original certification budget, and any shortcut that surfaces under the extended budget is reclassified as a shared blind spot rather than a genuine composition effect.

\subsection{Budget-Matched Adversarial Attacker}
\textbf{Challenge.} A single hand-constructed counterexample establishes that composition-induced shortcuts are possible, not how often they arise across the realistic combinations of tasks and verifier pairings a real pipeline might actually deploy --- and the space of pairings and shortcut strategies is too large to explore by constructing more examples by hand.

\textbf{Insight.} One fixed, identically applied search procedure, pointed at every single and composed gate across a benchmark and verifier-pairing grid, yields two distinct comparisons from the same underlying measurement: how the composed gate's unsound-acceptance behavior compares to its strictest single-verifier component, and how that same behavior compares to an alternative that spends the identical search effort hardening one verifier instead of adding a second.

\textbf{Method.} An identical attacker --- fixed model, fixed prompt scaffold, fixed attempt budget --- is applied to every single-verifier gate, every composed gate, and an adversarially hardened single-verifier alternative, across a grid of benchmark tasks and verifier-pairing types, with per-gate discovery curves logged over the attack budget. The resulting false-acceptance rates are compared pairwise within each verifier-pairing cell under matched budgets, both between composed and single gates and between hardened-single and naive-composed configurations.

\subsection{Composition-Gap Classifier}
\textbf{Challenge.} A shortcut that both verifiers would have missed independently, even without ever being composed, is a fundamentally different finding from one that arises specifically because composition displaced the solution space --- and this distinction persists even when the attacker's search effort is perfectly matched across gates, so it cannot be resolved by the re-attack check alone.

\textbf{Insight.} The same underlying data that defines each verifier's blind-spot region --- which mutants of the reference solution each verifier's checks can and cannot detect --- can be used both to classify why a discovered shortcut occurred and to compute, in advance of any adversarial search, how complementary two verifiers' blind spots are; the taxonomy label and the complementarity score are two views of one object, not two separate analyses.

\textbf{Method.} Every shortcut that survives the re-attack check is labeled with an ablation triple --- whether it passes verifier A's single gate, verifier B's single gate, and the composed gate --- and, if genuinely emergent, sorted into a small structural taxonomy (a weaker disjunctive branch admitting what the stronger branch would reject; a conjunctive constraint displacing solutions outside the region either single certificate covers; a facet-split where each verifier certifies a different property mistaken for full-task coverage; and a threshold-interaction class for score-based verifiers). A coverage-complementarity score, computed as one minus the Jaccard overlap of the two verifiers' mutant-kill sets over a shared mutant pool, is regressed against per-cell emergent-shortcut incidence, controlling for task difficulty and gate type.

\section{Experiments}

\subsection{Experimental Setup}

\textbf{Datasets.} A constructed bounded-DSL pilot suite (this paper's own artifact) is the only setting where per-verifier soundness is exhaustively decidable, and is therefore required for the existence question below, which is defined as a proof rather than a rate. HumanEval+ and MBPP+ (EvalPlus-amplified test suites) supply the $\geq 10^4$-input differential ground truth the scale-up certifier and attacker require, at the smallest per-task constraint complexity, forming the base cell of the scale-up grid. BigCodeBench contributes tasks that natively carry multiple constraints per task, giving the conjunctive-gate conditions a realistic multi-facet setting rather than an artificially bolted-on second check. A post-training-cutoff subset of LiveCodeBench isolates the prevalence and hardening-versus-adding findings from a specific rival explanation: that an apparent shortcut is really the attacker or verifier models reproducing a memorized near-solution rather than genuinely lacking a general one, since post-cutoff problems cannot have been memorized by the models involved. BigCodeBench and LiveCodeBench ship no EvalPlus-amplified corpus, so their differential ground truth is generated by applying the same EvalPlus-style amplification to each task's reference tests; any task whose corpus cannot reach the $\geq 10^4$-input floor is excluded and counted toward the thin-cell accounting below rather than silently retained.

\label{sec:baselines-placebo}\textbf{Baselines} (sanity-check comparators and rival-explanation controls, not competitors to beat --- this is a measurement study, and no research question here is framed as ``our procedure beats baseline X''). \emph{B0-C (current practice).} Verifier A alone and verifier B alone, each attacked on its own task variant at the identical budget, are the necessary ``why does this matter at all'' control: if composed gates behave indistinguishably from their strictest single component everywhere, the motivating concern collapses. \emph{B0-P (null-by-construction placebo, evaluated first).} Before any hypothesis is scored, the degenerate same-space configuration --- both single verifiers and the composed gate certified and attacked over one shared space (in the pilot) or a self-composition of two independently seeded instances of the same check (at scale) --- is run through the identical pipeline. Because a ground-truth-failing composed accept in this configuration would contradict a same-space single-gate certificate by construction (Problem Definition), this arm must return exactly zero emergent shortcuts; any nonzero count marks the run invalid (ground-truth incompleteness, a certificate bug, or an attacker bookkeeping error), and no hypothesis below is scored until the pipeline is re-diagnosed from the walking skeleton. Naive conjunctive and disjunctive gates are the compositions under direct test. An adversarially hardened single verifier --- instantiated training-free as iterative adversarial test augmentation (repeatedly attacking the current suite and adding a killing test for every discovered unsound acceptance, cost-matched to the naive-AND arm), approximating the objective of prior RL-trained verifier-hardening work \citep{utrl2508} without its training cost --- tests whether hardening one check is a substitute for adding a second. A thresholded, continuous-score reward model --- adapted from an execution-free reward-model design for software-engineering agents \citep{swerm2512} by fixing a pre-registered acceptance cutoff calibrated on a held-out split, not tuned against attack results, an explicit disclosed adaptation since the source model outputs a continuous ranking score rather than a Boolean verdict --- instantiates one concrete heterogeneous verifier pairing (a test suite paired with a thresholded learned scorer) in the scale-up grid, using a released checkpoint where available or else a prompted proxy, never a new training run. A structured, rule-linked aggregation scheme, adapted from a guardrail design that reasons explicitly over inter-check relationships rather than aggregating naively \citep{r2guard2407}, tests whether the composition gap comes from naive Boolean aggregation specifically, or from composition in general.

\textbf{Metrics.} The false-acceptance rate of a gate, $\mathrm{FAR}(G)$ = the probability that an attacker-found candidate is accepted by $G$ while failing the differential ground truth, at a stated attack budget; the emergent-gap rate, the fraction of (task, verifier-pairing) cells with at least one shortcut that passes the composed gate while both single gates hold at equal budget and the shortcut survives the $10\times$-budget re-attack; attack-success-at-budget, the probability the attacker finds at least one accepted-but-wrong candidate within the stated number of attempts; the coverage-complementarity score, one minus the Jaccard overlap of two verifiers' mutant-kill sets; and the certification budget at which each single-gate certificate was established. For stochastic (LLM-based) verifiers, a gate verdict is defined conservatively: an acceptance counts only if it is stable under three independent verdict draws at fixed decoding settings (3-of-3), per-verifier verdict-flip rates are reported, and any cell whose flip rate exceeds five percent is excluded from the primary comparisons and reported separately --- this is the verifier-determinism exclusion the thin-cell policy below refers to.

\textbf{Research questions:}
\begin{enumerate}[leftmargin=1.6em]
\item RQ1 --- Does a Boolean-composed gate admit a shortcut that neither single-verifier variant admits, over an exhaustively enumerated candidate space?
\item RQ2 --- Does re-attacking a shortcut's single-gate origin at ten times its certification budget reclassify some, but not most, emergent-labeled shortcuts as shared blind spots?
\item RQ3 --- Under matched attack budgets, does a composed gate show a higher false-acceptance rate than its strictest single-verifier component across a realistic benchmark and verifier-pairing grid?
\item RQ4 --- Is hardening a single verifier a substitute for adding a second, independently certified one?
\item RQ5 --- Does coverage complementarity between two verifiers predict where composition-induced gaps concentrate?
\item RQ6 --- Do disjunctive- and conjunctive-gate shortcuts concentrate in distinct, predictable structural classes?
\end{enumerate}

All six research questions are scored only after --- and conditional on --- the null-by-construction placebo B0-P (Baselines above) returning exactly zero emergent shortcuts; a nonzero placebo count invalidates the run before any question below is answered.

\subsubsection{RQ1: Does a Boolean-composed gate admit a shortcut that neither single-verifier variant admits, over an exhaustively enumerated candidate space?}
\textbf{Setup.} For each composition rule (conjunctive and disjunctive, scored independently), five hand-constructed pilot task families are exhaustively enumerated over their bounded candidate space. For each family and rule, the single-verifier-A and single-verifier-B variants are certified sound by full enumeration, and the corresponding composed variant is searched over the same fully enumerated space for a candidate the composed gate accepts and the ground truth rejects.

\textbf{Expected outcome.} For each composition rule scored independently, at least three of the five pilot task families contain at least one such shortcut, with both single-verifier variants certified shortcut-free over the identical fully enumerated space.

\textbf{Falsification condition.} Refuted for a rule if fewer than three of five families contain a certified shortcut for it; if refuted for AND only (or, mirrored, for OR only), the claim narrows to the composition rule that did produce a certified shortcut. If refuted for both rules \emph{and} RQ3 also shows no composed-versus-single effect, the finding pivots to the global fallback: composition preserved soundness in every constructed and benchmark condition, bounding rather than confirming the folklore composition-risk warning --- a pivot claimable only with its pre-registered credibility evidence (coverage reported over a parameterized task-family grid, discovery-versus-budget curves showing search saturation, and an explicit restriction of the null to the attacker class actually run), so that a negative result cannot be attributed to simply not looking hard enough. If refuted for both rules \emph{but} RQ3 still holds, the global fallback does \emph{not} fire (its pivoted claim would then be false): the pre-committed claim for this cell is instead that emergent shortcuts resisted construction in the bounded pilot DSL, yet a budget-matched adversary finds them at significantly higher rates on composed than on strictest-single gates across realistic benchmark stacks --- composition-induced unsoundness is a deployment-scale phenomenon whose existence evidence is the discovered benchmark shortcuts themselves, not pilot certificates, with the bounded-DSL negative reported as a limit of the construction method rather than of the phenomenon.

\textbf{Hypothesis (plain language).} The composition of two facet-certified, task-locally-sound verifiers is constructible to fail --- not merely conceivable in principle, but demonstrable as a machine-checked fact in at least a majority of deliberately constructed cases, for both composition rules. This is deliberately a claim about constructibility under the pre-registered extension rule (Problem Definition): its evidential role is to validate the instrument and fix the definitions, not to carry the paper's headline --- prevalence on realistic stacks, where no constructed extension exists, is RQ3's question.

\textbf{Metrics.} Binary per (task family, composition rule): existence of an accepted, ground-truth-failing candidate over the fully enumerated composed-variant space, alongside the corresponding single-verifier certificates.

\textbf{Result.} \ph

\subsubsection{RQ2: Does re-attacking a shortcut's single-gate origin at ten times its certification budget reclassify some, but not most, emergent-labeled shortcuts as shared blind spots?}
\textbf{Setup.} Every shortcut provisionally labeled emergent in the scale-up sweep (RQ3 below) has its single-verifier-A and single-verifier-B task variants re-attacked at ten times the original certification budget, using the same fixed attacker configuration.

\textbf{Expected outcome.} A nonzero but minority fraction --- under 50\% --- of provisionally emergent-labeled shortcuts are reclassified as shared blind spots under the extended-budget re-attack.

\textbf{Falsification condition.} Refuted (control is inert) if the reclassified fraction is 0\%; refuted (control explains away the finding) if the reclassified fraction is 50\% or more. Either result means the emergent-gap rate reported under RQ3 cannot be trusted at face value until the re-attack protocol itself is revisited.

\textbf{Hypothesis (plain language).} The re-attack control is doing real work --- it does catch some cases where an ``emergent'' label was really an under-searched single verifier --- without gutting the emergent-gap finding it exists to validate.

\textbf{Metrics.} Count and fraction of provisionally emergent-labeled shortcuts reclassified after the $10\times$ re-attack, reported per verifier-pairing cell.

\textbf{Result.} \ph

\subsubsection{RQ3: Under matched attack budgets, does a composed gate show a higher false-acceptance rate than its strictest single-verifier component across a realistic benchmark and verifier-pairing grid?}
\textbf{Setup.} A grid of approximately 200 benchmark tasks across the four datasets above, crossed with approximately six verifier-pairing types, is attacked with a fixed comparison attacker $\mathcal{A}$ (fixed model, fixed prompt scaffold, 500 attempts per gate) applied identically to every single and composed gate; single-gate certification uses the independent certification attacker $\mathcal{A}_{\text{cert}}$ (Problem Definition) so that the comparison attacks never draw randomness used to select what counts as certified-sound in the first place. Twelve primary paired comparisons (six pairing types $\times$ two composition rules) are tested via McNemar's exact test on discordant unsound-acceptance indicators, at $\alpha = 0.05$ with Holm--Bonferroni correction across the twelve. The power arithmetic behind the effect-size floor in the expected outcome below is pre-registered together with the assumption it depends on, and at the corrected threshold rather than the nominal one. For disjunctive-gate cells, discordance is structurally one-sided: every candidate found unsoundly accepted by a single gate is also evaluated against the corresponding composed gate, and monotonicity transfers a single-gate accept to the OR gate, so single-only unsound acceptances cannot occur; under one-sided discordance, McNemar's exact test needs at least nine discordant tasks to clear the most conservative Holm threshold ($\alpha/12$), which a seven-percentage-point composed-gate excess supplies with probability approximately 0.9 at 200 tasks per cell. For conjunctive-gate cells no such structural constraint holds --- single-only unsound acceptances (the single gate admits what the composed gate blocks) are expected, discordance is genuinely two-sided, and the same seven-point floor at the corrected threshold is powered at 0.8 only if the single-only discordance rate stays low (on the order of one percentage point); each cell's observed single-only rate is therefore reported, and the pre-registered response to a materially two-sided cell is to enlarge it toward roughly 340 tasks from the BigCodeBench pool, or, where that is infeasible, to report that cell's minimum detectable asymmetry at the corrected threshold instead of silently retaining the seven-point floor. \emph{Thin-cell policy.} If compute limits or verifier-determinism/attack-symmetry exclusions leave any of the twelve primary cells below 160 tasks (the smallest $n$ at which the one-sided-discordance arithmetic above still retains 0.8 power for a 7pp excess at the corrected threshold; conjunctive cells remain additionally subject to the observed single-only-rate caveat), that cell is pre-declared INCONCLUSIVE rather than scored as pass, fail, or silently pooled with an adjacent cell; the ``at least four of twelve'' count below is then evaluated only over non-INCONCLUSIVE cells, with the INCONCLUSIVE count reported alongside the headline. With $k < 12$ scorable cells, the pass threshold scales proportionally to at least $\lceil k/3 \rceil$ of the $k$ scorable cells (preserving the pre-registered four-of-twelve ratio), the requirement that at least one qualifying comparison is a conjunctive-gate comparison is unchanged, and a minimum of six scorable cells is required --- below that, this question itself is declared inconclusive rather than scored in either direction.

\textbf{Expected outcome.} In at least four of the twelve primary comparisons, including at least one conjunctive-gate comparison, the composed gate's false-acceptance rate exceeds the strictest single gate's by a discordant-pair asymmetry of at least seven percentage points, significant after correction.

\textbf{Falsification condition.} Refuted if fewer than four of the twelve comparisons clear the significance and effect-size bar, or if none of the qualifying comparisons is a conjunctive-gate comparison. If refuted while RQ1 still holds, the claim narrows to: composition-induced unsoundness is constructible and exhaustively certifiable, but a budget-matched adversary rarely finds it on realistic benchmark stacks --- a tail event conditioned on verifier-pair structure, not an ambient property of two-check pipelines.

\textbf{Hypothesis (plain language).} Composition-induced unsoundness is not confined to hand-built curiosities --- a budget-matched adversary finds it often enough, across a realistic grid, to register as a measurable rate rather than a rare anomaly, though not necessarily in every cell of the grid.

\textbf{Metrics.} Per-cell false-acceptance rate and 95\% confidence interval for each gate configuration; McNemar odds ratios with confidence intervals; count of significant cells after Holm correction.

\textbf{Result.} \ph

\subsubsection{RQ4: Is hardening a single verifier a substitute for adding a second, independently certified one?}
\textbf{Setup.} For each verifier-pairing type, an adversarially hardened single-verifier gate is attacked under the identical fixed-attacker configuration and compared, via a paired McNemar exact test at $\alpha = 0.05$ (Holm-corrected within this secondary comparison family), against the naive conjunctive gate of the corresponding unhardened pair, at equal total attack budget.

\textbf{Expected outcome.} For at least one verifier-pairing type, the hardened single gate's false-acceptance rate is at least five percentage points lower than the naive conjunctive gate's, significant after correction.

\textbf{Falsification condition.} Refuted if no verifier-pairing type shows the hardened-single gate beating the naive conjunctive gate by the stated margin at significance --- the claim would then become that adding a second facet-certified check never underperformed adversarially hardening the first, so the composition warning (RQ1/RQ3) concerns silent certificate invalidation, not a prescription to prefer single hardened gates.

\textbf{Hypothesis (plain language).} Spending a fixed amount of adversarial-hardening effort on one verifier and spending the same effort on adding a second, unhardened verifier are not interchangeable defenses --- and at least sometimes, hardening what you already have beats adding something new.

\textbf{Metrics.} Per-pairing-type false-acceptance rate for the hardened-single and naive-conjunctive configurations, with paired confidence intervals and test statistics.

\textbf{Result.} \ph

\subsubsection{RQ5: Does coverage complementarity between two verifiers predict where composition-induced gaps concentrate?}
\textbf{Setup.} For every (task, verifier-pairing) cell in the scale-up grid, a coverage-complementarity score is computed from the two verifiers' mutant-kill sets over a shared mutant pool. A logistic regression of emergent-gap incidence on this score, controlling for task difficulty and gate type, is fit across all cells, with a Wald test on the complementarity coefficient at $\alpha = 0.05$ (Holm-corrected within this secondary family); because the same tasks recur across pairing types, per-cell observations are not independent, and the Wald test uses cluster-robust standard errors clustered at the task level. Task difficulty is measured as one minus the fraction of fifty independent, gate-feedback-free samples from the certification attacker's base model that satisfy the ground truth, computed once per task and frozen before the sweep launches, so the covariate cannot absorb attack outcomes.

\textbf{Expected outcome.} The complementarity coefficient is positive and significant, with an odds ratio of at least 1.5 per one-standard-deviation increase in the complementarity score.

\textbf{Falsification condition.} Refuted if the complementarity coefficient is non-positive or fails significance after correction. If jointly refuted alongside RQ6's taxonomy-concentration test, the claim becomes: emergent shortcuts exist and occur at measurable rates, but are not predicted by mutant-kill-set complementarity nor concentrated by gate type --- no cheap pre-deployment warning signal was found, strengthening the case that composed gates must be attacked directly to be audited.

\textbf{Hypothesis (plain language).} Composition-induced gaps are not scattered unpredictably --- they concentrate where two verifiers are blind to different, non-overlapping regions of failure, which is measurable in advance of running any adversarial search at all.

\textbf{Metrics.} Logistic-regression coefficient and odds ratio for the complementarity score; Wald test statistic and corrected p-value.

\textbf{Result.} \ph

\subsubsection{RQ6: Do disjunctive- and conjunctive-gate shortcuts concentrate in distinct, predictable structural classes?}
\textbf{Setup.} Every genuinely emergent shortcut discovered in the scale-up sweep is assigned a structural taxonomy label (weak-branch, constraint-displacement, facet-split, or threshold-interaction) from its ablation-triple pattern and verifier-pairing configuration, double-annotated on a human-adjudicated sample with inter-annotator agreement required before scale-up labeling proceeds.

\textbf{Expected outcome.} At least 60\% of disjunctive-gate emergent shortcuts fall in the weak-branch class, and at least 50\% of conjunctive-gate emergent shortcuts fall in the constraint-displacement class, each significantly above the 25\% rate expected under a uniform distribution over the four classes.

\textbf{Falsification condition.} Refuted if either proportion fails to exceed the 25\% uniform-baseline rate at significance --- see RQ5 for the joint fallback claim if both this test and the complementarity regression fail together.

\textbf{Hypothesis (plain language).} Disjunctive and conjunctive composition fail in structurally different, nameable ways --- a weak fallback branch for disjunction, a constraint that pushes solutions somewhere uncertified for conjunction --- rather than failing through one generic, undifferentiated mechanism.

\textbf{Metrics.} Per-class proportion of emergent shortcuts for each composition rule, with exact binomial confidence intervals against the uniform baseline.

\textbf{Result.} \ph

\vspace{0.5em}
\noindent\textbf{Further evaluation.} With more time, a natural extension is testing pipelines with three or more heterogeneous checks and non-Boolean (weighted or thresholded) aggregation --- explicitly out of scope here beyond a single illustrative structured-aggregation control arm, because moving past two-verifier Boolean gates changes the combinatorics of the certification and attribution steps enough to need its own design rather than a direct extension of the protocol above. A second natural extension is auditing whether passive, checklist-style review (rather than active budget-matched search) can detect these gaps at all, since the design above only ever measures what an active adversary finds.

\section{Related Work}

\textbf{Composed reward and verifier systems already deployed.} A reward-chain framework \citep{rewardchain2601} converts a single verifiable answer into a hierarchy of content and style sub-rewards, aggregated by a simple weighted average, to avoid the hacking a single reward signal invites; a flawed-positive-aware policy optimizer \citep{fapo2510} adds a process-level detector to an outcome-only verifier to catch correct-answer-wrong-reasoning cases; an execution-free reward model \citep{swerm2512} is trained to complement execution-based test suites for software-engineering agents where tests are sparse or incomplete; and a reasoning-based guardrail \citep{r2guard2407} aggregates per-category safety classifiers, explicitly documenting that naive aggregation already under-estimates risk when unsafe intent is spread across several weak signals. All four validate their design by its downstream effect --- end-task performance gains, or a documented single-guardrail failure mode --- never by directly attacking the composed gate itself to test whether it, as a whole, remains sound; the guardrail work comes closest to this paper's concern but confines its fix to homogeneous categories inside one system, not the heterogeneous verifier pairs (tests, judges, execution-free scorers) this paper considers, and the reward-chain work's ``composition'' is continuous averaging rather than a hard Boolean gate at all --- both are treated here as method-twins requiring an explicit, disclosed adaptation (a pre-registered acceptance threshold, calibrated on a held-out split) before they can instantiate a Boolean verifier-B in our gate algebra, not as drop-in baselines.

\textbf{Reward-hacking theory for a single proxy.} A formal treatment of proxy-reward validity \citep{correlatedproxies2403} ties reward hacking to the breakdown of a proxy's correlation with the true objective under a reference policy, and shows that regularizing toward that reference policy provably bounds the resulting true-reward loss --- a single-proxy, single-reference-policy theory that never engages the discrete accept/reject algebra of a hard Boolean gate, and offers no account of what happens when two independently-valid proxies are combined.

\textbf{Hardening or mapping a single verifier's failure surface.} One line of work \citep{utrl2508} adversarially strengthens a test suite against the current solver's failure modes by rewarding discrimination between a reference implementation and sampled imperfect programs; another \citep{biasscope2602} automatically mines a judge's own blind-spot inventory from its misjudgments, iteratively perturbing rejected responses to surface biases beyond any predefined taxonomy. Both characterize one verifier's failure surface in isolation and in depth, but neither asks what happens when two independently mapped failure surfaces are combined into a single gate --- exactly the interaction this paper's structural taxonomy and complementarity score are built to characterize.

\textbf{Fundamental limits of a single verifier.} A theoretical analysis of scalable evaluation via LLM judges \citep{judgelimits2410} bounds the accuracy a proxy judge can certify by its correlation with ground truth, showing that correlation is necessarily bounded once the evaluated model outperforms the judge; a separate cryptographic result \citep{filterbarriers2507} proves that, under standard hardness assumptions, no external filter running faster than the model it screens can reliably distinguish an adversarially-constructed harmful prompt from a benign one. Both establish hard limits for a single verifier acting alone, but neither addresses how such a bound composes --- whether stacking a second, differently-limited check changes the barrier, tightens it, or opens a new one.

\textbf{Classical composition-failure literatures.} Outside the recent verification literature, older lines of work have studied failures of composed defenses: ensembles of weak defenses in adversarial machine learning, where composing individually circumventable defenses fails to compose robustness; N-version programming and fault-diversity reliability theory, where failure correlation between independently developed components undercuts the reliability gains that independence arguments promise; defense-in-depth failure analyses in security engineering; and the adaptive-attack evaluation discipline in adversarial machine learning, which insists a defense be scored only against attacks re-tuned to it at matched effort --- a methodological ancestor of the budget-matching protocol here. Closest of all in logical shape, compositional program verification internalized long ago that local correctness certificates do not compose for free: assume-guarantee and rely-guarantee reasoning exists precisely because a component's certificate holds only relative to an environment assumption, and compositions of individually verified components are unsound without explicitly discharged interface side-conditions. These literatures either compose components that are weak or only statistically characterized --- so their negative results concern aggregate robustness or reliability --- or, in the compositional-verification case, treat unsound composition as a proof obligation to be discharged inside one formal system whose components share an interface language. The setting here differs in exactly the property under test: each component enters composition carrying an individual, budget-stamped soundness certificate earned by search rather than proof, the components are heterogeneous checks (test suites, static analyzers, learned judges) that share no formal interface language in which an assume-guarantee side-condition could even be written down, and the estimand is the composed gate's certificate-relative false acceptance measured by matched-budget attack --- whether composition destroys a property each component demonstrably had, not whether it fails to add a property no component had alone. We flag these adjacencies explicitly rather than leaving them to be discovered in review; a systematic comparison against the strongest match in these literatures is future work.

\section{Conclusion}

This work turns a background assumption of modern verification pipelines --- that combining independently sound checks can only make a gate stricter --- into a question with a precise, testable structure, rather than an implicit default. The central commitment is that ``sound'' only means something once anchored to a specific facet, task variant, and budget, and that a composition failure counts as genuine only once unequal search effort and pre-existing shared blind spots are both actively ruled out. Whichever way the pre-registered comparisons resolve, the result is actionable: if composition reliably preserves soundness across the tested grid, the folklore fear attributed to multi-check pipelines is substantially bounded, at least for two-check Boolean gates; if it does not, the coverage-complementarity signal lets pipeline designers flag risky check pairings before deployment, not only after an adversary has found the gap. Four limitations bound this claim from the outset: the pilot's existence evidence is conditional on the pre-registered candidate-space extension rule and is accordingly treated as instrument validation rather than as the headline finding --- the claims that matter are the scale-up's rate and predictor, measured where no constructed extension exists; the design is scoped to two-verifier Boolean gates over checkable code tasks, leaving richer aggregation and non-code domains to future audits built on the same released instrument; every soundness claim outside the exhaustive pilot is a budget-relative statement, not a proof --- reported here, not resolved, and itself a future methodological target; and the hardened-single-verifier baseline (RQ4) is instantiated training-free, as iterative adversarial test augmentation rather than the RL-trained procedure of the prior work it approximates, which trades some fidelity to that trained baseline for keeping the design's ``no training runs required'' property intact --- a conceded approximation, not a claim that the two are identical.

\vspace{1em}
\hrule
\vspace{1em}

\section*{Pre-handoff checklist}
\begin{itemize}[leftmargin=1.4em]
\item \textbf{Organization/navigation} --- narrative flows problem $\to$ method $\to$ experiments $\to$ related work without forward references to undefined terms; the pilot/scale-up (exact vs. budget-relative) distinction is introduced once in Problem Definition and reused consistently, never re-derived mid-document.
\item \textbf{Terminology discipline} --- ``shortcut,'' ``facet-relative soundness,'' ``emergent,'' ``coverage complementarity,'' ``structural taxonomy'' each defined at first use in Problem Definition or Method Overview, used consistently thereafter; no internal stage IDs (\texttt{sp\_*}, \texttt{m\_*}, \texttt{rq\_*}, C1--C4, M1--M4, H1--H4) leaked into prose --- checked directly against the discussion pipeline log.
\item \textbf{Notation with worked example} --- the formal setup section defines $(T, r, D, GT, V, G, \mathrm{Sound}, \Pi, B_{\text{cert}})$ and the Problem Definition/Methodology sections consistently reuse this notation; no separate worked numeric example is included since no numbers exist yet (would require fabricating results).
\item \textbf{RQ framing matches outcomes} --- all 6 RQs are graded/quantitative (rates, proportions, odds ratios, reclassification fractions), each answered in its own subsection with an explicit falsifiable expected outcome, not a binary yes/no.
\item \textbf{No typesetting artifacts} --- only explicit placeholder markers in Results, per the drafting convention this is a pre-registered spec, not a report.
\item \textbf{Artifact plan} --- the bounded-DSL pilot suite, the task-pair generator, the certifier, and the attack logs are the natural release artifacts; not yet built, so not claimed as available.
\end{itemize}

\section*{Hand-off}
This is a draft, not a finished paper --- experiments have not run. Natural next steps: prose polish and venue-specific LaTeX formatting once a target venue is chosen, or revising an earlier stage if piloting surfaces a problem (e.g., if the bounded-DSL construction cannot simultaneously contain genuine solutions and shortcut-shaped candidates at a tractable enumeration size).

\bibliographystyle{plainnat}
\bibliography{refs}

\end{document}
